\begin{problem}[33届物理竞赛复赛]{61}
    光电倍增管是用来将光信号转化为电信号并加以放大的装置，其结构如图 a 所示：它主要由一个光阴极、n 个倍增级和一个阳极构成；光阴极与第 1 倍增级、各相邻倍增级及第 n 倍增级与阳极之间均有电势差 V；从光阴极逸出的电子称为光电子，其中大部分（百分比 $\eta$ ）被收集到第 1 倍增级上，余下的被直接收集到阳极上；每个被收集到第i倍增级（ $i=1,\cdots,n$ ）的电子在该电极上又使得 $\delta$ 个电子（ $\delta>1$ ）逸出；第i倍增级上逸出的电子有大部分（百分比 $\sigma$ ）被第 $i+1$ 倍增级收集，其他被阳极收集；直至所有电子被阳极收集，实现信号放大。已知电子电荷量绝对值为e。

    \begin{subproblem}
        求光电倍增管放大一个光电子的平均能耗，已知 $\delta\sigma>1,\quad n>>1$ 。
    \end{subproblem}
    \begin{subproblem}
        为使尽可能多的电子从第 $i$ 倍增级直接到达第 $i + 1$ 倍增级而非阳极，早期的光电倍增管中，会施加垂直于电子运行轨迹所在平面（纸面）的匀强磁场。设倍增级的长度为 $a$ 且相邻倍增级间的几何位置如图 b 所示，倍增级间电势差引起的电场很小可忽略。所施加的匀强磁场应取什么方向以及磁感应强度大小为多少时，才能使从第 i 倍增级垂直出射的能量为 $E_{e}$ 的电子中直接打到第 $i+1$ 倍增级的电子最多？磁感应强度大小为多少时，可以保证在第 $i+1$ 倍增级上至少收集到一些从第 i 倍增级垂直出射的能量为 $E_{e}$ 的电子？
    \end{subproblem}
\end{problem}